% Copyright (c) 2022-2026 MKM Research Labs. Licensed under the MIT License; see LICENSE.
%
% MKM-DD-001 Depth-Damage Curve Model -- model documentation v2.0, the document of record
% filed on the model as Core Documentation. It supersedes the filed v1.1 (18-Feb-2026),
% whose source reached v1.4 (26-Mar-2026, removed from PhysicalRisk with the other
% per-model docs in 527b5a348) without being filed. Rewritten after the Wolfram benchmark
% (MKM-DD-001_wolfram_benchmark v0.1) found the code had moved past both: v2.0 documents
% what src/models/floodrisk/ computes. Its sensitivity tables are inlined, generated
% from the code in the house generator's format.
%
% Build: it inputs the house header, which lives in PhysicalRisk. Copy
% PhysicalRisk/docs/models/shared beside this folder as ../shared, then run
% pdflatex three times (the third settles "Page n of N").
\documentclass[11pt,a4paper]{article}

\newcommand{\doctitle}{Depth-Damage Curve Model}
\newcommand{\docsubtitle}{Depth-Damage, Velocity, Attenuation and Spatial Functions of the Flood Risk Suite}
\newcommand{\docversion}{2.0}
\newcommand{\docdate}{03-November-2026}
\newcommand{\docauthor}{David K Kelly}

\input{../shared/mkm_header}
% T1 so the underscores in code identifiers print as underscores.
\usepackage[T1]{fontenc}
\usepackage{lmodern}
% The shared history table sets every column to its natural width; wrap the description.
\renewcommand{\mkmhistorypage}[1]{%
    \section{Document History}
    \begin{table}[ht]
        \centering
        \begin{tabular}{l|l|p{9cm}|l}
            \textbf{Date} & \textbf{Version} & \textbf{Description} & \textbf{Author} \\
            \hline
            #1
        \end{tabular}
    \end{table}
}

\begin{document}
% Long code paths and identifiers cannot hyphenate; let the spacing stretch instead.
\sloppy

\mkmtitlepage
\mkmlegalpage

\tableofcontents
\clearpage

% ===================================================================
\mkmexecsummary
% ===================================================================

This document presents the mathematical framework, calibration and validation of the
Depth-Damage Curve Model (MKM-DD-001), the functions in \texttt{src/models/floodrisk/}
that turn a flood at a gauge into damage at a property. It documents what the code
computes as of package version 2.1, and supersedes v1.1 of this document (filed as
\emph{Flood Risk Model}), which still described the property-type factor and the
distance-decayed depth calculation removed from the code in May 2026.

The model comprises:
\begin{enumerate}[nosep]
    \item \textbf{The depth-damage curve} --- a UK-calibrated piecewise-linear
          vulnerability curve from depth above floor level to damage ratio (\S\ref{sec:dd}).
    \item \textbf{The flood threshold} --- the property's height above its reference gauge
          plus its floor level (\S\ref{sec:threshold}).
    \item \textbf{The velocity functions} --- Manning's equation, ground slope, travel time,
          exponential retention with distance, and the single-peak property hydrograph
          (\S\ref{sec:velocity}).
    \item \textbf{The spatial functions} --- the haversine distance, inverse-distance
          weighting and an exponential correlation matrix (\S\ref{sec:spatial}).
    \item \textbf{Portfolio analytics} --- a correlated Monte Carlo of portfolio loss with
          VaR, ES, HHI and clustering (\S\ref{sec:portfolio}).
\end{enumerate}

\noindent The curve, the threshold and the velocity functions are what production uses:
the Property Flood Response model (MKM-PF-001) calls them for every property and storm,
and the hazard curve builder applies the curve to gauge levels. The spatial functions
other than the haversine, and the portfolio analytics, have no caller in production.

\textbf{Key sensitivity:} the damage curve scales expected loss one for one. A 10\%
shift of the curve moves a property's expected loss by exactly 10\% while its shocked
damage ratio cannot reach 1, which holds for damage ratios up to about 0.6; above that
the clip at 1 absorbs part of the shift (Table~\ref{tab:sens_shift}). The portfolio
figure therefore depends on the portfolio's mix of damage ratios.

% ===================================================================
\mkmpurpose
% ===================================================================

\subsection{Purpose}
The model quantifies the physical impact of fluvial flooding on a property: how deep the
water stands above its floor, when it arrives, and what fraction of reinstatement cost it
causes. Those damage ratios feed PRS pricing and the property flood timeseries.

\subsection{Scope}
\begin{itemize}[nosep]
    \item \textbf{In scope:} the depth-damage curve; the flood threshold; Manning velocity,
          slope, travel time and retention; the single-peak property hydrograph; the
          haversine, IDW and correlation functions; the portfolio Monte Carlo and its
          metrics.
    \item \textbf{Out of scope:} how a storm's gauge response is chosen and propagated to
          a property, including the bankfull offset and terrain scaling (MKM-PF-001); the
          compound per-pulse hydrograph (\texttt{floodrisk/hydrograph/}, MKM-HG-001); GEV
          hazard curve fitting; PRS pricing.
    \item \textbf{Upstream:} gauge responses and timeseries; property elevation and floor
          level.
    \item \textbf{Downstream:} MKM-PF-001 (property flood response); the hazard curve
          builder (\texttt{src/models/hazard/builder.py}, \texttt{terrain\_grid.py}).
\end{itemize}

% ===================================================================
\mkmframework
% ===================================================================

\subsection{Depth-Damage Curve}
\label{sec:dd}

The curve maps depth $d$ above the finished floor (metres) to a damage ratio
$\delta\in[0,1]$, the fraction of reinstatement cost incurred, through ten control points
calibrated against UK flood damage survey data:

\begin{table}[H]
\centering
\begin{tabular}{lcccccccccc}
    \toprule
    \textbf{Depth (m)} & 0 & 0.05 & 0.5 & 1.0 & 1.5 & 2.0 & 3.0 & 4.0 & 5.0 & 6.0 \\
    \midrule
    \textbf{Damage ratio} & 0 & 0.05 & 0.25 & 0.40 & 0.50 & 0.60 & 0.75 & 0.85 & 0.95 & 1.0 \\
    \bottomrule
\end{tabular}
\caption{UK-calibrated depth-damage control points}
\label{tab:dd_control}
\end{table}

Between adjacent control points $d_i\le d<d_{i+1}$,
\begin{equation}
    \delta(d) = \delta_i + \frac{d - d_i}{d_{i+1} - d_i} \, (\delta_{i+1} - \delta_i),
    \label{eq:dd}
\end{equation}
with $\delta(d)=0$ for $d\le0$ and $\delta(d)=1$ for $d\ge6$\,m.
\texttt{scalar\_depth\_damage} evaluates it with a plain loop over the control points
(\texttt{DEPTH\_POINTS}, \texttt{DAMAGE\_POINTS} in \texttt{config/damage.py}).

\paragraph{No property-type factor.} Earlier versions scaled the curve by 1.2 for
commercial and 0.9 for industrial property. The code no longer does; the factors remain
in config as \texttt{PROPERTY\_TYPE\_DAMAGE\_FACTORS} but nothing reads them.

\paragraph{The degree-5 curve.} \texttt{bri\_depth\_damage} evaluates a degree-5
polynomial through the origin, the least-squares fit to Table~\ref{tab:dd_control}
($R^2=0.999244$, RMSE 0.009368, largest error 0.020388 at 0.05\,m, increasing on
$[0,6]$), on the depth less the floor level and a signed BRI stilt. It is implemented for
the BRI models and has no caller in production; the curve of record is
Eq.~\eqref{eq:dd}.

\subsection{Flood Threshold}
\label{sec:threshold}

A property floods when water at it rises above its height above the reference gauge plus
its floor level (\texttt{relative\_elevation}):
\begin{equation}
    h_{\mathrm{thr}} = \max(0,\ e_{\mathrm{prop}} - e_{\mathrm{gauge}}) + f,
    \label{eq:threshold}
\end{equation}
with ground elevations $e$ in metres AOD and floor level $f$ in metres. The depth at the
property is the water above that threshold, and the damage ratio is
$\delta(\max(0,\ \text{water} - h_{\mathrm{thr}}))$.

\paragraph{Removed in v2.0.} Versions 1.x derived depth from the flood WSE, the property's
elevation and a linear decay over 25\,km from a Thames reference point (their
Eqs.~5--6), capped at 5\,m. The code replaced that with the nearest-gauge propagation of
MKM-PF-001 and Eq.~\eqref{eq:threshold}; \texttt{FLOOD\_MAX\_DISTANCE\_M} remains in
config, unread.

\subsection{Velocity, Retention and the Property Hydrograph}
\label{sec:velocity}

\subsubsection{Manning's Equation}
For wide, shallow overbank flow the hydraulic radius is approximated by the depth $h$:
\begin{equation}
    v = \frac{1}{n}\, h^{2/3} S^{1/2},
    \label{eq:manning}
\end{equation}
with $v=0$ when $h\le0$ or $n\le0$, and the slope taken as $\max(|S|,\,0.001)$.

\subsubsection{Ground Slope}
\begin{equation}
    S = \max\!\left(\frac{|e_{\mathrm{gauge}} - e_{\mathrm{prop}}|}{d},\ 0.001\right),
    \label{eq:slope}
\end{equation}
and $S=0.001$ when $d\le0$.

\subsubsection{Travel Time}
\begin{equation}
    t = \frac{d}{3600\, v}\ \text{hours},
    \label{eq:travel}
\end{equation}
infinite when $v=0$.

\subsubsection{Retention with Distance}
The share of the flood signal that reaches a property at distance $d$ from the river:
\begin{equation}
    a(d) = e^{-d/\ell},\qquad \ell = 10{,}000\ \text{m},
    \label{eq:retention}
\end{equation}
with $a=1$ for $d\le0$ and $a=0$ for $\ell\le0$. At 600\,m a property retains 94\% of the
signal, at 2\,km 82\%, at 5\,km 61\%. There is no near-field bypass: v1.4's rule that
$a=1$ within 2\,km has been removed. (v1.1 gave $\ell=2{,}000$\,m.)

\subsubsection{Property Hydrograph}
\texttt{build\_property\_hydrograph} turns the gauge's hourly series into the property's.
With gauge base $W_b$ (its minimum), peak $W_p$ at hour $t_p$, travel time $t_r$ and the
property's peak WSE $W^{*}$:
\begin{itemize}[nosep]
\item before the flood front arrives, $t<t_r$, the property sits at $W_b$ with no depth;
\item on the rising limb, $t\le t_p+t_r$, the gauge is sampled at the nearest hour to
      $t-t_r$;
\item on the recession, at the nearest hour to $t_p+(t-t_p-t_r)/\gamma$, $\gamma=1.5$;
\end{itemize}
and the property's level and depth are
\begin{equation}
    W(t) = W_b + (W^{*} - W_b)\, \frac{W_{\mathrm{gauge}}(\cdot) - W_b}{W_p - W_b},\qquad
    d(t) = \max\bigl(0,\ W(t) - e_{\mathrm{prop}} - f\bigr),
    \label{eq:hydrograph}
\end{equation}
both reported to the millimetre. Retention is \emph{not} applied here: the caller
(MKM-PF-001) applies $a(d)$ to the water above bankfull before computing $W^{*}$, so the
hydrograph scales by it once. v1.x's Eq.~15 multiplied by $a$ again.

\subsection{Spatial Functions}
\label{sec:spatial}

\subsubsection{Haversine Distance}
\begin{equation}
    d = 2R\,\operatorname{atan2}\!\left(\sqrt{a},\ \sqrt{1-a}\right),\quad
    a = \sin^2\frac{\Delta\varphi}{2} + \cos\varphi_1\cos\varphi_2\sin^2\frac{\Delta\lambda}{2},
    \label{eq:haversine}
\end{equation}
with $R=6{,}371{,}000$\,m. This is the one spatial function production calls; it is
used across the property, gauge and seismic code.

\subsubsection{Inverse-Distance Weighting}
Given gauges at distances $d_i$ with water surface elevations $W_i$,
$\hat W=\sum_i w_iW_i/\sum_iw_i$ with $w_i=d_i^{-p}$, $p=2$; the first gauge closer than
1\,m returns its own $W_i$. Production replaced IDW with the single nearest gauge in
package v2.1; the function is retained.

\subsubsection{Correlation Matrix}
For properties at (longitude, latitude) in degrees, with Euclidean separation
$d^{\deg}_{ij}$ converted at 111\,km per degree,
\begin{equation}
    \rho_{ij} = \rho_0 \exp\!\left(-\frac{111\, d^{\deg}_{ij}}{L/1000}\right),\quad i\ne j,
    \qquad \rho_{ii}=1.
    \label{eq:corr}
\end{equation}
The degree conversion takes no account of latitude: at 51.5$^\circ$N it overstates
east--west separation by 61\%. The matrix is positive definite.

\subsection{Portfolio Analytics}
\label{sec:portfolio}

\texttt{risk\_analytics} simulates portfolio loss under correlated shocks. It has no
caller in production; it is documented because it is part of the package.

\paragraph{Monte Carlo.} With damage ratios $\delta_j$, values $V_j$ and correlation
matrix $\Sigma$, it draws $M=1{,}000$ shock vectors $z\sim N(0,\Sigma)$ with NumPy's
legacy generator (seed 42), which factorises $\Sigma$ by singular-value decomposition,
and sets
\begin{equation}
    L^{(m)} = \sum_j V_j\, \mathrm{clip}\bigl(\delta_j(1+0.2z_j^{(m)}),\,0,\,1\bigr).
    \label{eq:mc}
\end{equation}
VaR$_\alpha$ is the $\alpha$ percentile of the draws (NumPy's linear interpolation) and
ES$_\alpha$ the mean of the draws at or above it, for $\alpha=95,99$.

\paragraph{The expected loss, exactly.} Each shocked ratio $X_j=\delta_j(1+0.2z_j)$ is
normal with mean $\delta_j$ and deviation $0.2\delta_j$, so
\begin{equation}
    \mathbb{E}\,\mathrm{clip}(X_j,0,1) = \delta_j\bigl[\Phi(\beta_j)-\Phi(-5)\bigr]
    + 0.2\delta_j\bigl[\varphi(-5)-\varphi(\beta_j)\bigr] + 1-\Phi(\beta_j),\quad
    \beta_j = \frac{1-\delta_j}{0.2\,\delta_j},
    \label{eq:el}
\end{equation}
and the expected portfolio loss is $\sum_jV_j\,\mathbb{E}\,\mathrm{clip}(X_j,0,1)$ whatever
the correlation. The Monte Carlo mean estimates it.

\paragraph{Concentration.} The Herfindahl--Hirschman index of losses
$\ell_j=V_j\delta_j$ is $\sum_js_j^2$ with $s_j=\ell_j/\sum_k\ell_k$; it is $1/N$ for an
equal portfolio. Grid concentration aggregates by 1\,km cells (1$^\circ\approx$111\,km),
and KMeans ($K=5$, ten initialisations, seed 42) clusters properties on standardised
depth, damage, position and $\ln(1+V)$.

% ===================================================================
\mkminputs
% ===================================================================

\subsection{Model Parameters}
\begin{table}[H]
\centering
\small
\begin{tabular}{llll}
    \toprule
    \textbf{Parameter} & \textbf{Symbol} & \textbf{Default} & \textbf{Where} \\
    \midrule
    Manning's roughness & $n$ & 0.04 & \texttt{DEFAULT\_ROUGHNESS} \\
    Minimum slope & $S_{\min}$ & 0.001 & \texttt{MIN\_SLOPE} \\
    Retention length & $\ell$ & 10{,}000\,m & \texttt{DEFAULT\_RETENTION\_LENGTH} \\
    Recession factor & $\gamma$ & 1.5 & \texttt{DEFAULT\_RECESSION\_FACTOR} \\
    Earth radius & $R$ & 6{,}371{,}000\,m & \texttt{haversine\_distance} \\
    IDW power, minimum distance & $p$, $d_{\min}$ & 2, 1\,m & \texttt{idw\_interpolate\_from\_gauges} \\
    Correlation length, base & $L$, $\rho_0$ & caller's & \texttt{build\_correlation\_matrix} \\
    Monte Carlo draws, seed & $M$ & 1{,}000, 42 & \texttt{simulate\_portfolio\_impact} \\
    Shock dispersion & --- & 0.2 & \texttt{simulate\_portfolio\_impact} \\
    KMeans clusters, restarts & $K$ & 5, 10 & \texttt{analyze\_risk\_clusters} \\
    \bottomrule
\end{tabular}
\caption{Model parameters (\texttt{config/models/\_flood.py}, \texttt{config/damage.py})}
\label{tab:params}
\end{table}

\subsection{Depth-Damage Calibration}
The control points of Table~\ref{tab:dd_control} follow the Multi-Coloured Manual
(Penning-Rowsell et al., 2013), chosen to capture minimal damage below 0.05\,m
(below-floor wetting), rapid escalation from 0.05 to 1.0\,m (electrics, ground-floor
contents, plaster), and saturation above 3\,m as the structure approaches total loss.

\subsection{Manning's Roughness}
The default $n=0.04$ represents a suburban floodplain with mixed land cover (Chow, 1959):
smooth concrete 0.012, short grass 0.025, urban and suburban 0.035--0.050, dense
vegetation 0.070--0.100.

\subsection{Retention Length}
$\ell=10{,}000$\,m was chosen so that properties within a few kilometres of the river
keep most of the flood signal (82\% at 2\,km), with the terrain scaling of MKM-PF-001
shortening or lengthening the effective distance by land cover. It is an expert
judgement, not fitted to observed inundation extents.

% ===================================================================
\mkmimplementation
% ===================================================================

\subsection{Module Architecture}
\begin{table}[H]
\centering
\small
\begin{tabular}{lp{95mm}}
    \toprule
    \textbf{Module} & \textbf{Responsibility} \\
    \midrule
    \texttt{depth\_damage.py} & \texttt{scalar\_depth\_damage} (Eq.~\ref{eq:dd}); the BRI
        polynomial, stilt and floor uplift; \texttt{is\_prs\_flood} \\
    \texttt{elevation.py} & \texttt{relative\_elevation} (Eq.~\ref{eq:threshold}) \\
    \texttt{velocity.py} & Manning, slope, travel time, retention and the property
        hydrograph (Eqs.~\ref{eq:manning}--\ref{eq:hydrograph}) \\
    \texttt{spatial.py} & haversine, IDW, correlation matrix, KDTree, nearest point on
        a polyline \\
    \texttt{risk\_analytics.py} & Monte Carlo, VaR, ES, HHI, grid concentration, KMeans \\
    \texttt{hydrograph/} & the compound per-pulse hydrograph (MKM-HG-001; out of scope) \\
    \bottomrule
\end{tabular}
\caption{Modules of \texttt{src/models/floodrisk/} (package version 2.1)}
\label{tab:modules}
\end{table}

\subsection{How Production Uses It}
For each property and storm, MKM-PF-001 takes the controlling gauge's peak, subtracts
bankfull, scales by the retention $a(d)$ of Eq.~\eqref{eq:retention}, compares the
result with the threshold of Eq.~\eqref{eq:threshold}, computes the travel time from
Eqs.~\eqref{eq:manning}--\eqref{eq:travel}, builds the hydrograph of
Eq.~\eqref{eq:hydrograph}, and applies the curve of Eq.~\eqref{eq:dd} to the peak depth.
The hazard curve builder applies Eq.~\eqref{eq:dd} to the gauge level above the severe
threshold.

\subsection{Numerical Considerations}
\begin{itemize}[nosep]
    \item \textbf{Slope floor:} $S_{\min}=0.001$ keeps Manning's velocity finite on flat ground.
    \item \textbf{Zero velocity:} the travel time is infinite; MKM-PF-001 then uses 0.
    \item \textbf{Damage clipping:} shocked ratios are clipped to $[0,1]$.
    \item \textbf{Shock draws:} the legacy \texttt{np.random.multivariate\_normal} factorises
          by SVD, which tolerates a semi-definite matrix.
    \item \textbf{Reproducibility:} the Monte Carlo reseeds NumPy's global generator with
          42 on every call.
\end{itemize}

% ===================================================================
\mkmsensitivity
% ===================================================================

Tables~\ref{tab:sens_dd}--\ref{tab:sens_ret} are generated from the code; each value is
rounded and then printed by size (two places from 1 up, four below).

\begin{table}[H]
    \centering
    \caption{Depth-damage function: damage ratio by flood depth.}
    \label{tab:sens_dd}
    \begin{tabular}{rr}
        \toprule
        \textbf{Depth (m)} & \textbf{Damage ratio} \\
        \midrule
        0.0000 & 0.0000 \\ 0.0500 & 0.0500 \\ 0.1000 & 0.0722 \\ 0.2500 & 0.1389 \\
        0.5000 & 0.2500 \\ 0.7500 & 0.3250 \\ 1.00 & 0.4000 \\ 1.50 & 0.5000 \\
        2.00 & 0.6000 \\ 3.00 & 0.7500 \\ 4.00 & 0.8500 \\ 5.00 & 0.9500 \\ 6.00 & 1.00 \\
        \bottomrule
    \end{tabular}
\end{table}

\begin{table}[H]
    \centering
    \caption{Manning's velocity (m/s) by roughness coefficient ($S=0.005$).}
    \label{tab:sens_manning}
    \begin{tabular}{rrrr}
        \toprule
        \textbf{Manning's $n$} & \textbf{$v$ at $d=0.5$m} & \textbf{$v$ at $d=1.0$m} & \textbf{$v$ at $d=2.0$m} \\
        \midrule
        0.0200 & 2.23 & 3.54 & 5.61 \\
        0.0300 & 1.49 & 2.36 & 3.74 \\
        0.0400 & 1.11 & 1.77 & 2.81 \\
        0.0500 & 0.8910 & 1.41 & 2.25 \\
        0.0600 & 0.7420 & 1.18 & 1.87 \\
        0.0800 & 0.5570 & 0.8840 & 1.40 \\
        \bottomrule
    \end{tabular}
\end{table}

\begin{table}[H]
    \centering
    \caption{Retention factor by distance and characteristic length.}
    \label{tab:sens_ret}
    \begin{tabular}{rrrrr}
        \toprule
        \textbf{Distance (m)} & \textbf{$\ell=3000$m} & \textbf{$\ell=5000$m} & \textbf{$\ell=10000$m} & \textbf{$\ell=20000$m} \\
        \midrule
        0 & 1.00 & 1.00 & 1.00 & 1.00 \\
        500 & 0.8465 & 0.9048 & 0.9512 & 0.9753 \\
        1000 & 0.7165 & 0.8187 & 0.9048 & 0.9512 \\
        2000 & 0.5134 & 0.6703 & 0.8187 & 0.9048 \\
        5000 & 0.1889 & 0.3679 & 0.6065 & 0.7788 \\
        10000 & 0.0357 & 0.1353 & 0.3679 & 0.6065 \\
        \bottomrule
    \end{tabular}
\end{table}

\subsection{Depth-Damage Curve Shift}
Scaling every damage ratio by $1\pm10\%$ and taking the exact expected loss of
Eq.~\eqref{eq:el} gives, for a single property:

\begin{table}[H]
    \centering
    \caption{Expected loss after a $\pm$10\% curve shift, relative to the base curve.}
    \label{tab:sens_shift}
    \begin{tabular}{rrr}
        \toprule
        \textbf{Damage ratio $\delta$} & \textbf{+10\%} & \textbf{$-$10\%} \\
        \midrule
        0.1--0.5 & 1.1000 & 0.9000 \\
        0.6 & 1.0997 & 0.9000 \\
        0.7 & 1.0947 & 0.9009 \\
        0.8 & 1.0785 & 0.9074 \\
        0.9 & 1.0559 & 0.9228 \\
        \bottomrule
    \end{tabular}
\end{table}

\noindent Below $\delta\approx0.6$ the shocked ratio almost never reaches 1 and the loss
moves exactly with the curve. A portfolio's response is the value-weighted mix of these
rows; v1.x's ``$\pm$8\%'' corresponds to a portfolio dominated by damage ratios near 0.8.

\subsection{Manning's Roughness}
At depth 1\,m and slope 0.003, over a 3\,km flow path:

\begin{table}[H]
    \centering
    \caption{Velocity and travel time by roughness ($d=1$\,m, $S=0.003$, 3\,km).}
    \label{tab:sens_rough}
    \begin{tabular}{rrr}
        \toprule
        \textbf{$n$} & \textbf{Velocity (m/s)} & \textbf{Travel time (h)} \\
        \midrule
        0.025 & 2.19 & 0.38 \\
        0.040 & 1.37 & 0.61 \\
        0.060 & 0.9130 & 0.91 \\
        0.100 & 0.5480 & 1.52 \\
        \bottomrule
    \end{tabular}
\end{table}

\noindent Travel time scales with $n$. v1.x printed velocities of 0.64, 0.40, 0.27 and
0.16\,m/s for this case, which Eq.~\eqref{eq:manning} does not give.

\subsection{Retention Length}
At 2\,km from the river a property retains 0.37 of the signal with $\ell=2{,}000$\,m,
0.67 with 5{,}000\,m and 0.82 with the default 10{,}000\,m. Halving $\ell$ squares the
retention at a given distance, so the flooded footprint is sensitive to it; that is why
$\ell$ is a governed parameter.

\subsection{Portfolio Sensitivities}
The correlation length and Monte Carlo sample size sensitivities of v1.x (its Tables~11
and~12) were computed on a portfolio the document did not specify, and are withdrawn. The
Monte Carlo mean estimates Eq.~\eqref{eq:el}, so its sampling error can be judged against
that exactly for any portfolio.

% ===================================================================
\mkmvalidation
% ===================================================================

\subsection{Depth-Damage Curve Validation}
As reported in v1.x: the curve lies within the Multi-Coloured Manual's 25th--75th
percentile band for residential property at every depth; against Environment Agency
post-event surveys (summer 2007, winter 2013/14) its mean absolute error is 6\% in
damage ratio over 340 properties; and 5-fold cross-validation on the MCM data gives an
RMSE of 0.07.

\subsection{Velocity Validation}
As reported in v1.x: against TUFLOW 2D outputs on three Thames floodplain transects,
Manning's velocity is within 18\% on average, overestimating in complex urban terrain;
travel times are within $\pm$30 minutes over 2\,km. The overestimate is conservative,
bringing the flood front earlier.

\subsection{Implementation Benchmark}
The Wolfram oracle benchmark (\texttt{MKM-DD-001\_wolfram\_benchmark}) verifies the code
against Eqs.~\eqref{eq:dd}--\eqref{eq:el} at 50 digits and through an independent Wolfram
Engine evaluation, and reproduces every printed figure in
Tables~\ref{tab:sens_dd}--\ref{tab:sens_rough}.

% ===================================================================
\mkmlimitations
% ===================================================================

\begin{enumerate}
    \item \textbf{Fluvial only.} Surface water, groundwater, coastal and sewer flooding
          are not modelled.
    \item \textbf{Static curve.} Damage depends on depth alone, not on duration,
          velocity, contamination or warning time.
    \item \textbf{Piecewise-linear curve.} It may understate damage between 0.05 and
          0.5\,m, where escalation is nonlinear.
    \item \textbf{Shallow-flow velocity.} $R_h\approx h$ holds only for wide, shallow
          overbank flow.
    \item \textbf{Isotropic retention.} Eq.~\eqref{eq:retention} decays with distance
          alone, ignoring preferential flow paths; $\ell$ is expert judgement.
    \item \textbf{Gaussian dependence.} The Monte Carlo's normal shocks are symmetric and
          may understate joint tail losses; the 0.2 dispersion is uncalibrated.
    \item \textbf{Degree conversion.} Eq.~\eqref{eq:corr} and the grid concentration use
          1$^\circ\approx$111\,km in both directions, overstating east--west distance by
          61\% at 51.5$^\circ$N.
    \item \textbf{No flood defences} and \textbf{single events only}.
    \item \textbf{Unused code.} IDW, the KDTree, the correlation matrix, the portfolio
          analytics and the BRI polynomial have no production caller, and the property-type
          factors and 25\,km decay radius sit unread in config.
\end{enumerate}

% ===================================================================
\mkmgovernance
% ===================================================================

\subsection{Q1 --- Ronseal Test}
Yes, for what production uses: the curve, the threshold and the velocity functions turn
a flood at a gauge into depth, timing and damage at a property. The portfolio analytics
do what they say but are not used.

\subsection{Q2 --- Accuracy and Stability}
The code computes its documented formulae to rounding error (\S7.3); the curve is
validated against MCM and EA data and the velocity against TUFLOW (\S7.1--7.2). The
functions are deterministic; the Monte Carlo is seeded.

\subsection{Q3 --- Resource Usage}
The production functions are scalar and cheap. The correlation matrix is $N\times N$:
800\,MB for 10{,}000 properties.

\subsection{Q4 --- Override Controls}
No manual overrides. Parameters are version-controlled constants in \texttt{config/};
changes go through code review.

\subsection{Q5 --- Output Consumption}
Damage ratios and flood depths feed the property flood timeseries and PRS pricing
through MKM-PF-001, and the hazard curves through the hazard builder.

\subsection{Q6 --- Model Chaining}
Yes: MKM-PF-001 and the hazard curve builder consume it; it consumes gauge responses
and property elevations.

\subsection{Q7 --- Weaknesses}
See \S8; the most material are the static curve, the expert-judged retention length
and fluvial-only coverage.

\subsection{Q8 --- Monitoring}
Annual review of the curve against EA and MCM data; review of $\ell$ against observed
inundation extents when available; re-running the implementation benchmark on any
change to \texttt{src/models/floodrisk/}.

\subsection{Q9 --- Overall Risk}
Acceptable, conditional on the review of $\ell$: the curve is the most material
component and is validated, while the retention length directly sets which properties
flood and is not yet fitted to data.

% ===================================================================
\mkmchangelog
% ===================================================================

\begin{itemize}[nosep]
    \item \textbf{v2.0:} documents package v2.1: the property-type factor and the
          distance-decayed depth removed; the flood threshold of
          Eq.~\eqref{eq:threshold}; $\ell=10{,}000$\,m with no near-field bypass;
          retention applied once, by the caller; SVD shocks; the exact expected loss;
          the curve-shift, roughness and retention tables regenerated from the code;
          v1.x's correlation and sample-size tables withdrawn; the degree conversion
          error corrected to 61\%.
\end{itemize}

\mkmhistorypage{13-Feb-2026 & 1.0 & Initial documentation & D. Kelly \\
18-Feb-2026 & 1.1 & Added Section 9: Governance Validation Questions & D. Kelly \\
26-Feb-2026 & 1.2 & Delta Engine added to inventory; models in scope updated (not filed) & D. Kelly \\
26-Mar-2026 & 1.4 & Nearest-gauge propagation; retention length 10 km with a 2 km near-field bypass (not filed) & D. Kelly \\
03-Nov-2026 & 2.0 & Rewritten to the code (package v2.1) after the Wolfram benchmark & David K Kelly \\}

\end{document}
